% ===========================================================================
\section{Background and Motivation}
\label{sec:background}
% ===========================================================================

\subsection{Slow Faults in Distributed Systems}

Slow faults represent a spectrum of performance degradations that fall between
normal operation and complete failure~\cite{huang-gray-failure}. Unlike crash
faults, which trigger immediate failover, slow faults can persist undetected
while progressively destabilizing system invariants. Prior work has identified
several categories~\cite{xinda,gunawi-cbs}:

\para{Network slow faults} Additional latency on network links due to
congestion, switch buffer overflow, or NIC firmware issues. Even microsecond-scale
delays can disrupt heartbeat-based failure detectors when they cause cumulative
timeout expiration in consensus protocols~\cite{raft,zab}. Cloud providers
report tail latency spikes of 1--10\,ms occurring multiple times per day on
inter-AZ links~\cite{dean-tail-latency}.

\para{Filesystem slow faults} Increased I/O latency from disk degradation, file
system journaling contention, or storage controller issues. These particularly
affect write-ahead log (WAL) performance in databases. Lu~\etal{} found that
storage-related performance bugs account for 25\% of cloud system
issues~\cite{lu-cloud-reliability}. Modern NVMe SSDs can exhibit periodic
latency spikes of 10--100$\times$ normal due to garbage collection and wear
leveling~\cite{dean-tail-latency}.

\para{CPU slow faults} Resource throttling from cgroup limits, noisy
neighbors in multi-tenant environments, or garbage collection pauses that
starve critical threads. Kubernetes CPU throttling affects over 60\% of
production containers according to industry surveys~\cite{datadog-monitoring},
creating unpredictable scheduling delays for latency-sensitive consensus
heartbeats.

\para{Memory pressure faults} Reduced available memory forcing increased
garbage collection in JVM-based systems (Cassandra, HBase, Kafka, ZooKeeper),
page cache eviction increasing disk I/O, or swap thrashing. These create
cascading effects where a memory-constrained node becomes slow, triggering
failure detection by peers.

\para{Process faults} Container or process restarts and stops that cause
temporary unavailability. While these are detectable, the \emph{timing}
of restart relative to in-flight operations determines whether data loss
or corruption occurs.

\subsection{Xinda: Exhaustive Grid Exploration}
\label{sec:xinda-background}

\xinda{}~\cite{xinda} (NSDI~'25) provides the most comprehensive framework for
automated slow-fault testing. Its \emph{danger-zone} exploration uses a fixed
severity grid to systematically evaluate system behavior under increasing fault
intensity. The pipeline:

\begin{enumerate}[leftmargin=*,topsep=2pt,itemsep=1pt]
  \item Deploys the target system in Docker containers using pre-built images.
  \item Injects faults at each grid severity using Blockade (\texttt{tc netem})
        for network or CharybdeFS (FUSE) for filesystem.
  \item Executes a benchmark workload (YCSB, system-specific, or custom).
  \item Collects logs and metrics during execution.
  \item Reports which grid points trigger observable failures.
\end{enumerate}

The network danger-zone grid consists of 37 discrete values from
\texttt{slow-100$\mu$s} to \texttt{slow-1s}:

\begin{center}
\small
\texttt{100$\mu$s, 200$\mu$s, \ldots, 900$\mu$s, 1ms, 2ms, \ldots, 9ms,
10ms, 20ms, \ldots, 90ms, 100ms, 200ms, \ldots, 1s}
\end{center}

This approach is thorough for systems whose boundaries fall within the grid,
but has structural limitations:

\para{Fixed range} The grid covers exactly 100\,$\mu$s to 1\,s. Systems with
boundaries above 1\,s (Redis at 2.7\,s) or below 100\,$\mu$s produce
uninformative results---either all-miss or all-hit across the entire grid.

\para{Fixed granularity} Grid spacing is non-uniform (100\,$\mu$s steps at the
low end, 100\,ms steps at the high end). The actual boundary resolution depends
on where in the grid the boundary falls.

\para{No directionality} \xinda{} tests all 37 points regardless of results.
If the first point reproduces, the remaining 36 trials provide diminishing
information.

\para{System-agnostic grid} The same grid is applied to all systems. This is
efficient for systems with boundaries in the 1--100\,ms range but wasteful for
both very sensitive systems (boundary $\ll$ 100\,$\mu$s) and very tolerant
systems (boundary $\gg$ 1\,s).

\subsection{The Anduril Approach}

\anduril{}~\cite{anduril} (ATC~'25) takes a complementary approach: given a
known bug (typically a crash/exception bug with a code-level fix), it uses
feedback from runtime traces to guide fault injection toward reproducing the
specific bug. \anduril{} constructs a directed acyclic graph of causal events
and preferentially injects faults at points most likely to trigger the target
failure path.

While effective for exception bugs with discrete triggers, \anduril{}'s
approach requires code-level instrumentation (function interception, return
value modification) and targets single-event triggers rather than the
\emph{sustained} performance degradation characteristic of slow faults. A
slow fault manifests over many operations---there is no single function call
to intercept.

\subsection{Problem Statement}
\label{sec:problem}

Given:
\begin{itemize}[leftmargin=*,topsep=2pt,itemsep=1pt]
  \item A distributed system $S$ deployed in containers.
  \item A symptom oracle $O$ that evaluates logs and returns
        \emph{reproduced}/\emph{not reproduced}.
  \item A fault parameter space $\mathcal{F}$ with dimensions: type, location,
        severity, duration, start time.
\end{itemize}

Find: A fault configuration $f^* \in \mathcal{F}$ such that:
\begin{enumerate}[leftmargin=*,topsep=2pt,itemsep=1pt]
  \item $O(S, f^*) = \text{reproduced}$ (correctness).
  \item For each dimension $d$ of $f^*$, reducing $d$ causes
        $O(S, f') = \text{not reproduced}$ (local minimality).
  \item The total number of trial executions is minimized (efficiency).
\end{enumerate}

\noindent This formulation produces a \emph{locally minimal} recipe---minimal
along each dimension independently. Finding the \emph{globally minimal} recipe
(Pareto-optimal across all dimensions simultaneously) would require exponential
trials; we show that local minimality is sufficient for practical use.

\subsection{Motivating Observations}

Three observations motivate \system{}'s design:

\para{Observation 1: Fault effects are monotone} For additive faults (network
delay via \texttt{tc netem}, filesystem latency via injection), the effect is
monotonically non-decreasing with severity. If injecting 20\,ms delay causes
a leader election, then 25\,ms, 50\,ms, and 100\,ms also cause it---the higher
delay strictly subsumes the lower. This \emph{monotonicity property} enables
binary search over continuous severity ranges.

We formalize this as: for fault parameter $s$ and oracle $O$, if $O(s_1) =
\text{reproduced}$ and $s_2 > s_1$, then $O(s_2) = \text{reproduced}$. We
validate this assumption empirically across all 10 systems in
Section~\ref{sec:monotonicity}.

\para{Observation 2: Precision adapts to budget} A binary search with $k$
iterations over range $[0, S_{\max}]$ produces precision $S_{\max}/2^k$. With
8 iterations from 5000\,ms, precision is 19.5\,ms. With 12 iterations, it is
1.2\,ms. Operators can choose their precision--time trade-off explicitly,
unlike fixed grids where precision is predetermined by grid design.

\para{Observation 3: Minimal recipes are more actionable} A recipe stating
``19.5\,ms network delay for 30\,s triggers leader election'' is directly
useful for: (a)~CI regression tests that verify the boundary hasn't shifted,
(b)~SLA definitions that set alerting below the boundary, and (c)~architectural
decisions about timeout configuration. The recipe format is a concrete,
executable specification.
